\chapter{化学的计数单位：一“摩尔”有多少}

\begin{kaichang}
去厨房拿三样东西：一袋米、一小勺盐、一杯水。

先问你一个简单的问题：这袋米有多少粒？你多半会皱眉——谁有工夫一粒一粒数？但如果换个问法：这袋米有多重？往厨房秤上一放，十秒钟就有答案。

再问你一个不可能的问题：这杯水里有多少个水分子？这次连皱眉都不够用了。就算把全世界的人全动员起来，每人每秒数一个分子，不吃不睡数上一万年，也数不到这杯水的零头。

可是化学家偏偏能回答第二个问题，而且工具就是一台秤。他们还发明了一个专门的计数单位，名字叫“摩尔”。先猜一个悬念留着：一杯水里的分子个数，和全世界海洋能装下的“杯水”总数，哪个更多？这一章结束时，我们用数字对答案。
\end{kaichang}

\section{数数这件事，人类早就会偷懒}

你其实早就用过“打包计数”了。买鸡蛋论“打”：一打是 12 个；买牛奶论“箱”：一箱 24 盒；打印纸论“包”：一包 500 张。没人说“给我来 0.83 打鸡蛋”，因为“打”这个单位存在的意义，就是让数数变快、让算账变省事。这样的“打包”在生活中随处可见：纸巾论“抽”、邮票论“版”、电池论“板”，加油站干脆不数“滴”，直接按升卖油——升，其实也是液体世界里的一种大包。

再想想米。超市卖米从来不数粒，只称重。卖钉子的小店更绝：你买一盒钉子，店员往秤上一倒——他不是懒，他是在用\textbf{称量代替计数}。这个方法能成立，只需要知道一件事：一个钉子平均多重。整堆质量除以单个质量，个数就出来了，一颗都不用数。

\begin{shiyan}
\textbf{材料}：一台厨房秤（能精确到 1 克更好）、一小碗大米、一张纸、一支笔。

\textbf{步骤}：
\begin{enumerate}[itemsep=1pt]
\item 数出 100 粒米，放在纸上，称出总质量，记下来（比如 2.1 克）。
\item 算一算：平均每粒米多少克？
\item 抓一把米放在秤上，称出质量（比如 37 克），先别数，用刚才的“每粒质量”推算这把米有多少粒，把预测写在纸上。
\item 再真的数一数这把米，和预测比一比。
\end{enumerate}

\textbf{观察点}：预测值和真实值接近吗？差了多少粒？误差从哪里来——是不是有的米饱满、有的干瘪？如果你用 500 粒米来做第 1 步，预测会不会更准？

\textbf{安全提示}：这个实验没有任何危险，唯一的风险是数米数到眼花。做完把米倒回米缸，别浪费粮食。
\end{shiyan}

刚才你做的事情，和化学家每天做的事情，在思路上一模一样：\textbf{用一个“单个质量”当汇率，在“总质量”和“个数”之间兑换}。区别只在于，化学家面对的不是米粒，而是比米粒小得离谱的东西。

\section{原子小到什么程度}

离谱到什么程度？拿一个碳原子来说：它的质量大约是 $1.99\times 10^{-23}$ 克。注意那个指数：小数点后面要先写 22 个 0，才轮到 199。世界上最灵敏的分析天平，能称出约 $10^{-7}$ 克的质量——这已经是一台汽车价格的仪器了，可离一个原子还差着 16 个数量级。\textbf{单个原子是称不出来的，永远称不出来}。

再拿一粒食盐来说。一粒盐大约重 $6\times 10^{-5}$ 克（60 微克），小到肉眼几乎看不清。可它里面有多少个粒子？盐由钠离子和氯离子一对一对整齐堆成（第 18 章会细讲这种排列），后面你会算出来：这一小粒盐里大约有 $10^{18}$ 对离子——一百亿亿对。如果你每秒数 3 个，日夜不停地数，数完需要一百多亿年，比宇宙的年龄还长。

水分子更夸张。一杯水按 200 克算，里面有多少个水分子？先把答案的架子搭在这里：大约是 $7\times 10^{24}$ 个——七千万亿亿。全世界的人每人每秒数一个，要数三万年。

再换一个角度感受这种“小”：可见光的波长是原子直径的几千倍，原子比光波的“波纹”还细，所以无论显微镜镜片磨得多好，光都会像海浪绕过小石子一样绕过原子——原子不是“暂时还没被看见”，而是用光原则上就看不见。看不见、数不动、称不出单个，三件事把化学家逼上了同一条路。

所以化学家面对一个死局：反应明明是一个粒子一个粒子发生的（两个氢分子配一个氧分子，不多不少），可粒子小到没法数、没法称单个。怎么办？米粒实验已经给出提示：\textbf{不数，改称；只是需要一个巨大的“汇率”}。这个汇率就是本章的主角。

\section{碳原子当秤砣：相对原子质量}

要搭起“称量”和“计数”之间的桥，先得给每种原子报一个“体重”。可是原子的真实体重全是 $10^{-23}$ 克这种没法念的数，用起来太痛苦。于是化学家干脆换个思路：\textbf{不比绝对体重，比相对体重}——选一种原子当标准秤砣，别的原子和它比。

被选中的秤砣是碳-12 原子（碳原子中最常见的一种，原子核里 6 个质子、6 个中子；第 7 章会细讲同位素这件事）。规定：一个碳-12 原子的相对质量\textbf{恰好等于 12}。别的原子跟它比：氢原子的相对质量约为 1，氧约为 16，铁约为 56。这些数就叫\textbf{相对原子质量}，在周期表每个格子里都印着。

为什么要费这个劲？奥妙在下一步。一个碳-12 原子的实际质量是 $1.99\times 10^{-23}$ 克。现在问：多少个碳-12 原子堆在一起，正好重 12 克？做一次除法：
\[
\frac{12}{1.99\times 10^{-23}}\approx 6.02\times 10^{23}
\]
答案大约是 $6.02\times 10^{23}$ 个。这个数在化学里的地位，堪比圆周率在几何里的地位，它有自己的名字：\textbf{阿伏伽德罗常数}，符号 $N_{\mathrm{A}}$。2019 年起，国际单位制干脆把它钉死为一个精确值：
\[
N_{\mathrm{A}}=6.02214076\times 10^{23}\ \ud{}{mol^{-1}}
\]
（末尾的 $\mathrm{mol^{-1}}$ 读作“每摩尔”，马上解释。）

有了这个数，化学的计数单位就诞生了：\textbf{$6.02\times 10^{23}$ 个粒子，叫 1 摩尔}。摩尔是“物质的量”的单位，就像“个”是数量的单位、“打”也是数量的单位一样——只不过这一“包”有 6020 亿亿亿个。

\section{摩尔像“打”，但有三点不同}

把摩尔理解成“打”的放大版，方向是对的，但有三个关键差别：

\textbf{第一，数量级完全不是一个世界}。一打是 12，一摩尔是 $6.02\times 10^{23}$。这个数大到什么程度？后面有个估算会专门震撼你一次。

\textbf{第二，“打”是约定俗成，摩尔是精确定义}。12 这个数没有为什么，古代商人顺手定的；而 $6.02214076\times 10^{23}$ 是被国际单位制钉死的精确值，它锚定在碳-12 上（锚定的方式在 2019 年做了一次大修，见章末挑战层）。

\textbf{第三，也是最妙的一点}：因为阿伏伽德罗常数本身就是从“12 克碳-12”里算出来的，所以一个极其好用的性质自动成立——\textbf{1 摩尔任何粒子，用克称出来的质量，数值上恰好等于这种粒子的相对质量}。氢原子相对质量是 1，1 摩尔氢原子就重 1 克；氧原子相对质量 16，1 摩尔氧原子重 16 克；一个水分子相对质量是 $1\times 2+16=18$，1 摩尔水分子就重 18 克。分子也一样用：把化学式里各原子的相对质量加起来，就是这种分子的相对分子质量，也是它 1 摩尔的克数。

这就是化学家“用秤数原子”的全部秘密。想知道一杯 200 克的水里有多少个水分子？不用数：$200\div 18\approx 11.1$ 摩尔，再乘上 $6.02\times 10^{23}$，约 $6.7\times 10^{24}$ 个。开场那个不可能的问题，两行除法就解决了。

\begin{qiaoliang}
质量、物质的量、粒子数，是三间相邻的屋子，两座桥各有一个“汇率”：
\[
n=\frac{m}{M}\qquad N=n\times N_{\mathrm{A}}
\]
其中 $m$ 是质量（克），$M$ 是摩尔质量（每摩尔的克数，数值等于相对质量），$n$ 是物质的量（摩尔），$N$ 是粒子数，$N_{\mathrm{A}}$ 是阿伏伽德罗常数。

练两个来回。9 克水是多少摩尔？$M(\chem{H_2O})=18\ \mathrm{g/mol}$，所以 $n=9\div 18=0.5$ 摩尔，即约 $3.01\times 10^{23}$ 个分子。反过来，3 摩尔二氧化碳（$M(\chem{CO_2})=12+16\times 2=44\ \mathrm{g/mol}$）重多少？$3\times 44=132$ 克。

每算完一道题，做两个检查。\textbf{量纲检查}：$m\div M$ 的单位是 $\mathrm{g}\div(\mathrm{g/mol})=\mathrm{mol}$，单位对上了公式才没记反；要是算出 $\mathrm{g^2/mol}$ 这种怪物，立刻回去找错。\textbf{数量级检查}：日常烧杯里的样品是克级的，算出来的 $n$ 应该在 $10^{-2}$ 到 $10^{1}$ 摩尔之间；如果你算出 $10^{10}$ 摩尔，那比全世界的水还多，肯定是哪里错了。这两个检查不花十秒钟，能拦下考试里一半的低级错误。
\end{qiaoliang}

\section{三个类比，感受这个数有多大}

数字 $6.02\times 10^{23}$ 写在纸上平平无奇，可它大得完全超出日常直觉。用三个类比校准一下你的“数量级雷达”。

\textbf{沙子}。假设一粒沙是直径约 0.5 毫米的小球，体积约 $6.5\times 10^{-11}$ 立方米。1 摩尔粒沙子堆在一起，体积约 $4\times 10^{13}$ 立方米。把它们均匀铺在中国约 960 万平方千米的国土上，能铺出大约 4 米厚——整个国家埋在一层楼深的沙子里，这就是 1 摩尔粒沙。

\textbf{时间}。1 摩尔秒有多长？约 $1.9\times 10^{16}$ 年。宇宙的年龄约 138 亿年，也就是 $1.38\times 10^{10}$ 年。1 摩尔秒等于一百多万个宇宙年龄。换句话说，从宇宙大爆炸那一刻起每秒数一个数，数到今天，连一摩尔的零头都数不到。

\textbf{水}。1 摩尔水只有 18 克，一口就能喝完——这是“小”的方向；而一杯 200 克的水里却装着约 $7\times 10^{24}$ 个分子——这是“大”的方向。摩尔这个单位妙就妙在，它恰好站在两个方向的中间：一边连着你能喝、能称的宏观世界，一边连着十的二十几次方的粒子世界。

正因为单个原子小得令人绝望，才需要一个大得令人绝望的“包”。这两个绝望在 $6.02\times 10^{23}$ 这里恰好抵消，剩下的 18 克水、2 克氢气、56 克铁，全都是台秤上顺手的日常量。这不是巧合——这正是当初把标准定在“12 克碳-12”的原因。

\section{符号登场，再请出“浓度”}

到现在为止，本章的符号其实只有四个，一起亮个相：物质的量 $n$（单位 mol），阿伏伽德罗常数 $N_{\mathrm{A}}$，摩尔质量 $M$（单位 $\mathrm{g/mol}$），粒子数 $N$。它们之间的关系用一张“兑换地图”就能画全：

\begin{center}
\begin{tikzpicture}
\node[draw,rounded corners] (m) at (0,0) {质量 $m$（克）};
\node[draw,rounded corners] (n) at (5.4,0) {物质的量 $n$（摩尔）};
\node[draw,rounded corners] (N) at (11.6,0) {粒子数 $N$（个）};
\draw[->,thick] (m) -- (n) node[midway,above] {$\div M$};
\draw[->,thick] (n) -- (N) node[midway,above] {$\times N_{\mathrm{A}}$};
\draw[->,thick,dashed] (n.south) to[bend right=20] node[midway,below] {$\times M$} (m.south);
\draw[->,thick,dashed] (N.south) to[bend right=20] node[midway,below] {$\div N_{\mathrm{A}}$} (n.south);
\end{tikzpicture}
\end{center}

这只是人为的表示法：方框和箭头是帮助记忆的地图，不是物质的任何真实图像。

化学家还经常需要回答另一类问题：\textbf{多少东西，装在多少空间或多少混合物里}——这就是浓度。你已经见过它的好几种说法：生理盐水标签上的“$0.9\%$”，意思是每 100 克溶液里有 0.9 克氯化钠；空气预报里的“每立方米多少微克”，是质量除以体积；化学实验室最常用的是\textbf{摩尔浓度}：1 升溶液里溶了多少摩尔的溶质，记作 $\mathrm{mol/L}$。算一下就知道，生理盐水每升含 9 克氯化钠，$M(\chem{NaCl})\approx 58.5\ \mathrm{g/mol}$，所以浓度约为 $9\div 58.5\approx 0.154\ \mathrm{mol/L}$。

浓度还有一副“稀释”的面孔。把一勺浓汤倒进一锅水里，味道立刻变淡——溶质的量没变，体积变大了，浓度就被稀释了。实验室和农田都靠这个原理干活：农药瓶上写着“稀释 1000 倍使用”，意思是取 1 份药液兑约 1000 份水，把浓度降到既能杀虫、又尽量不伤作物、不在土壤里过度残留的区间。饮料瓶的配料表也在报浓度：一罐可乐大约每 100 毫升含 10 克糖，一罐 330 毫升的可乐里就有 30 多克糖，约合 0.1 摩尔——接近 20 块方糖。浓度这个数字本身没有好坏之分，剂量才决定一切：同样是氯化钠，$0.9\%$ 的溶液可以输进血管，而浓盐水喝下去只会让你脱水。

浓度用摩尔来报数，在医学上天天救场。化验单上的空腹血糖大约是 $5.6\ \mathrm{mmol/L}$（mmol 是毫摩尔，千分之一摩尔）。一个成年人全身血液约 5 升，算一算：血糖总量约 $5.6\times 5=28$ 毫摩尔，葡萄糖的摩尔质量是 $180\ \mathrm{g/mol}$，折合成质量只有约 5 克——大约一块方糖。你全身的血糖，居然只有一块方糖那么多，而身体要把它时刻维持在这个狭窄的范围内（第 60 章讲稳态时会回来讲这台精密的调节机器）。如果不换算成摩尔，你根本看不出“$5.6\ \mathrm{mmol/L}$”和“一块方糖”是同一回事。

\section{一个被冷落了半个世纪的假说}

\begin{zhengju}
阿伏伽德罗常数冠着阿伏伽德罗的名字，但他本人既没测过这个数，也不知道它有多大。1811 年，这位意大利物理学家提出一个假说：同温同压下，相同体积的气体含有相同数目的分子。这个假说能漂亮地解释气体反应的体积关系，可它要求承认“同种原子能结合成分子”（比如两个氢原子组成一个氢分子）——这与当时权威的观点冲突，于是假说被晾在一边将近五十年。原子论本身的来龙去脉，第 6 章会细讲，这里只说关键转折：1860 年，各国化学家在德国卡尔斯鲁厄开会，吵成一团，连“水的化学式是 HO 还是 \chem{H_2O}”都没有共识。会上意大利化学家坎尼扎罗散发了一本小册子，用阿伏伽德罗的假说把原子量的混乱账理顺了。许多年轻化学家读后当场“倒戈”——统一的相对原子质量标准这才可能建立。

那 $6.02\times 10^{23}$ 这个数本身是怎么测出来的？1865 年，奥地利人洛施密特用气体分子运动论估算了空气分子的数目，得了第一笔粗账。但真正让人信服的，是法国物理学家佩兰 1908 年前后的工作。他把藤黄树脂磨成的微小颗粒悬浮在水里，在显微镜下逐颗观察这些颗粒的布朗运动（第 2 章讲过这种永不停歇的乱舞）。佩兰发现，颗粒在重力下并不是均匀分布的，而是越往上越稀少，而且稀少的规律，和大气的密度随高度衰减的规律一模一样——仿佛这些颗粒是一群“放大了的原子”。用这条规律，他可以反推出：要让颗粒跳得这么欢，撞击它们的水分子得有多少个。他算出的阿伏伽德罗常数在 $6\times 10^{23}$ 到 $7\times 10^{23}$ 之间。

这个结论之所以致命，是因为佩兰用了好几种互不相干的方法——布朗运动、沉降平衡、放射性计数——条条大路都通向同一个数。原本坚决否认原子真实存在的科学家（比如主张“原子只是记账工具”的实证主义者），面对这种“不同方法互相印证”的阵势也先后认输。佩兰因此获得 1926 年诺贝尔物理学奖。这个故事最该记住的：一个看不见的东西是否真实存在，不靠辩论，靠的是\textbf{多条互相独立的证据链是否收敛到同一个答案}。顺便一提，佩兰的实验装置极其朴素：一台普通显微镜、一个加高的水柱、一根测微尺，加上近乎无穷的耐心——他逐颗数了几千个颗粒的高度。伟大的测量不一定需要昂贵的仪器，但一定需要让怀疑者挑不出毛病的设计。
\end{zhengju}

\section{这座桥的边界}

任何好用的模型都有边界，“摩尔”也不例外。

第一，\textbf{摩尔要求你数的东西是明确的、同种的粒子}。1 摩尔氢原子、1 摩尔水分子，都有意义；但“1 摩尔沙子”没有意义——沙子里混着成分不同、大小各异的颗粒，没有统一的“单个”可言。“1 摩尔淀粉”同样尴尬：淀粉分子是由成百上千个葡萄糖单元连成的大分子，每条链的长度都不一样，没有固定的相对分子质量（第 43、46 章还会遇到这类高分子）。

第二，\textbf{周期表上的相对原子质量大多是平均值}。氯的相对原子质量印着 35.5，并不是说存在一种“35.5 号”的氯原子，而是天然氯气里两种同位素按比例混合的平均结果——单个氯原子要么是约 35，要么是约 37。这个细节第 7 章会展开，现在只要记住：平均值拿来算账完全够用，但它不是任何一个真实原子的体重。

第三，\textbf{摩尔只负责计数，不负责别的}。1 摩尔不同物质，粒子数相同，但质量不同、体积也多半不同（气体在相同温度和压强下是个例外，那正是阿伏伽德罗假说的内容，第 4 章的气体图景和第 26 章的定量计算还会用到它）。

这些边界并不推翻摩尔，只是提醒你：任何单位都是为某类问题而生的。环保部门报空气污染用“微克每立方米”，关心的是总质量；毒理学家却常常要把它换算成“每立方米多少个分子”，因为造成伤害的是分子与身体的一次次碰撞，次数才和个数直接相关。同一份污染，两种报数，回答的是不同的问题——选对计数方式，本身就是一种科学素养。

\begin{wuqu}
“一摩尔任何物质都一样重。”——这是对摩尔最常见的误会。1 摩尔氢气（\chem{H_2}）只有 2 克，1 摩尔黄金却有 197 克，差了近百倍。摩尔保证的是\textbf{粒子个数相同}：1 摩尔氢气和 1 摩尔黄金，都含有 $6.02\times 10^{23}$ 个基本粒子（前者是分子，后者是原子）。个数相同不等于质量相同，就像“一打鸡蛋”和“一打西瓜”个数相同，你却不会指望它们一样重。同理，“摩尔”既不是质量单位，也不是体积单位——它是个计数单位，只此一家。
\end{wuqu}

\section{回到米袋、盐和那杯水}

现在兑现开头的承诺。

那袋米：你其实已经在用“摩尔思想”了。称出 100 粒米的质量，你就知道了“每粒米的平均质量”，整袋质量一除，粒数就出来了。化学家做的是同一件事，只是他们的“单粒质量”小到 $10^{-23}$ 克量级，所以兑换汇率大到 $10^{23}$ 量级。\textbf{阿伏伽德罗常数，本质上就是“原子世界的每粒米质量”的倒数}。

那粒盐：重约 $6\times 10^{-5}$ 克，除以氯化钠的摩尔质量 $58.5\ \mathrm{g/mol}$，约含 $10^{-6}$ 摩尔、也就是约 $10^{18}$ 对钠离子和氯离子。数是数不完的，但用摩尔说，它不过是“百万分之一摩尔”。

那杯水：200 克水约含 $6.7\times 10^{24}$ 个分子。回到开场的悬念——全世界的海洋约 $1.4\times 10^{24}$ 克水，按每杯 200 克算，只能装出约 $7\times 10^{21}$ 杯。\textbf{一杯水里的分子数，比全世界海洋能装下的杯水总数还要多，多大约一千倍}。换句话说，你把一杯水倒进大海，搅拌均匀后，再从海里任意舀起一杯水，里面平均能找到上千个原来那杯水的分子。你现在喝的这口水里，多半就有被恐龙喝过、被李白举杯邀过的分子——这不是诗，是除法。原子的世界就是这样：数量大到让“永恒循环”变成一句平淡的统计事实。

\section{为什么化学家非数个数不可}

你可能会问：既然称质量这么方便，为什么还要费劲绕到“个数”上去？因为\textbf{化学反应是粒子与粒子之间一对一、二对一地发生的}。两个氢分子恰好烧掉一个氧分子，一个碳原子恰好结合两个氧原子——反应的“配方”是按个数写的，不是按质量写的。拿合成氨厂来说：反应是 1 个氮分子配 3 个氢分子，如果工人按“1 克氮气配 3 克氢气”来配料，那就全错了——氮气的摩尔质量是 28，氢气只有 2，正确的配法是每 28 克氮气配 6 克氢气。不按个数换算，贵的原料会剩下一大半白白浪费，或者反应不完残留在产品里。药厂配药、污水处理厂加药，道理一模一样：都要先把质量换算成摩尔，配料比例才有意义。医院里也一样：输液的电解质浓度、透析液的配方、血糖和血钾的化验值，全用摩尔或毫摩尔报数，因为身体里的化学反应同样按“个数”记账。

那么，这种“按个数的账本”具体怎么记？一个化学方程式的系数到底在说什么？反应物会不会有剩余、剩余多少？这些问题的答案，就是摩尔思想真正的用武之地——第 26 章《化学方程式是反应的账本》，我们到时候把这本账彻底算清楚。

\begin{tiaozhan}
2019 年 5 月 20 日，国际单位制做了一次大手术：阿伏伽德罗常数被规定为\textbf{精确值} $6.02214076\times 10^{23}\ \mathrm{mol^{-1}}$，不再测量、永不修正。这听起来只是定义游戏，背后却是一次“主客易位”：在此之前，千克的定义是一个放在巴黎的铂铱合金圆柱体（“国际千克原器”），而摩尔是“12 克碳-12 所含的原子数”——也就是说，摩尔曾经依附于千克。可原器会沾灰、会磨损，一百多年里它和复制品之间出现了几十微克的漂移，用一块会变的金属定义全世界的质量，终究不是办法。2019 年之后，局面反过来：千克改由基本物理常数定义，摩尔由固定数目的粒子定义，铂铱圆柱体光荣退休。这场变革里最漂亮的工程之一，是科学家用高纯度硅-28 拉成单晶、磨成近乎完美的一千克球体，用 X 射线测出晶格间距，一格一格地“数”出球里有多少个硅原子——这是人类最接近“真的数了一遍原子”的实验。跳过本节完全不影响主线，但如果你好奇“定义”和“测量”谁说了算，这里有一座桥通向物理学。
\end{tiaozhan}

\section*{练一练}

\begin{sikaoti}
\item 照抄并补全本章的“兑换地图”（质量、物质的量、粒子数三个方框和四座箭头桥），然后在每座桥旁写一句“过桥费由谁收”（即需要知道哪个量才能换算）。这张图以后每一章定量计算都用得上。
\item 不查资料，估算一滴水（约 0.05 克）里有多少个水分子。再算：如果用一台每秒能数 100 万个的超级机器来数这些分子，要数多少年？
\item 4 克氢气（\chem{H_2}）和 32 克氧气（\chem{O_2}），哪种气体的分子更多？分别含有多少个分子？（提示：先求摩尔质量，分子个数比的是 $n$，不是质量。）
\item 有位同学记错了公式，写成 $n=M\div m$。不用代入数字，只做量纲检查，证明他一定错了。
\item 把 9 克氯化钠（\chem{NaCl}）溶解后加水到总体积 1 升，摩尔浓度是多少？和化验用的生理盐水（$0.9\%$）比一比，你发现了什么？
\item 估算“1 摩尔粒大米”的质量（一粒米约 0.02 克），再和全世界一年的粮食总产量（约 $3\times 10^{12}$ 千克）比较。这个对比能让你体会阿伏伽德罗常数为什么只能属于原子世界。
\end{sikaoti}
